\subsection{A construction-bound main theorem}
\label{ex-construction-headline}

The following theorem retains the connection between classical observables,
lattice representatives and quantum fields. Its normalizations and local
statements are specified below. Theorem~\ref{thm:ex-construction-assembly}
assembles their proofs on one retained Wilson family and one limiting
hierarchy.

Let $V_d$ be the finite-dimensional space of gauge-invariant polynomials
in radial curvature jets of weighted degree at most $d$, a curvature
derivative of order $j$ having weight $j+2$; see
Sections~\ref{sec:local-splitting} and~\ref{sec:local-sections}.
Fix a real polynomial basis, homogeneous in the number of curvature factors,
and use its complexification as coordinates on $V_d$, with the constant
polynomial as the unit coordinate.
A local section $L_d$ assigns to $v\in V_d$ a polynomial in the entries of
finitely many links and their conjugates, linearly in $v$, invariant under
lattice gauge transformations. Its stencil is independent of spacing and
period. It preserves the unit, conjugation and signed coordinate symmetries.
For every radial jet $J$ it satisfies the normalization
\begin{equation}\label{eq:headline-jet-normalization}
 \left.\frac{d^k}{dh^k}L_d(v;H(J,h))\right|_{h=0}
 =\left.\frac{d^k}{dh^k}v(h\cdot J)\right|_{h=0},\qquad 0\le k\le d.
\end{equation}
Here $H(J,h)$ is parallel transport along lattice links for the polynomial
radial connection at scale $h$, and $h\cdot J$ multiplies the weight-$j+2$
component by $h^{j+2}$. Thus small-field testing recovers the prescribed
classical polynomial.

Write $a_N=D_N^{-1}$ for the lattice spacing, $P_p$ for the physical
period, and $U$ for a link configuration on the corresponding torus.
For each $d$ choose a reference level $r_d$. At cutoff $(N,p)$, $N\ge r_d$,
the renormalized insertion has the form
\begin{equation}\label{eq:headline-insertion}
 I^{(d)}_{N,p}(v;x;U)=a_N^{-d}
 L_d\bigl(D_{\theta_N^{-1}}R_{d,N,p}^{-1}v;x;U\bigr).
\end{equation}
Here $\theta_N>0$ is the small-field threshold and $D_t$ multiplies a
coordinate with $s$ curvature factors by $t^s$. The invertible linear map
$R_{d,N,p}$ is fixed by transporting the inserted Wilson density from
level $N$ to $r_d$ and reading its polynomial coefficient there. This is
exact transport of the unnormalized inserted density by the same block
integrations as the uninserted density, with the scale factors and remainder
bounds of Sections~\ref{sec:local-splitting},
\ref{sec:terminal-marked-construction} and~\ref{sec:ex-marked-field-local}.
In particular, transporting $R_{d,N,p}^{-1}v$ gives coefficient $v$ at the
reference level. The maps in \eqref{eq:headline-insertion} must have this
transport and normalization; freely chosen renormalization matrices are
not an alternative definition. The same maps are used at every translated
anchor. Choices at the finitely many levels $N<r_d$ do not affect the limits.

\begin{remark}[The local normalization has content]
\label{rem:headline-nonconstant-sections}
If $L_d(v)$ is constant, with value $c$, then
\eqref{eq:headline-jet-normalization} says that the polynomial
$h\mapsto v(h\cdot J)$ has constant term $c$ and all other coefficients
zero: its degree is at most $d$. Taking $h=1$ gives $v(J)=c$ for every
$J$. In particular $L_d$ is injective, and a nonconstant curvature
polynomial cannot have a constant local representative. Since both $D_t$
for $t>0$ and $R_{d,N,p}$ are invertible, the insertion space in
\eqref{eq:headline-insertion} contains nonconstant functions whenever
$V_d$ contains nonconstant polynomials. On sufficiently large tori the
stencil does not wrap, so its link assignments extend to torus
configurations and this remains true there. This rules out the specific
replacement of \emph{all} lattice insertions by constants. It does not
by itself prove that the continuum field space contains a field independent
of the unit, or that its correlations are non-Gaussian.
\end{remark}

\begin{theorem}[The Wilson construction]
\label{thm:ex-construction-headline}
For every compact matrix group $G$ with nonabelian simple Lie algebra as
specified above, and every positive blocking penalty $b$, one can choose
one sequence of lattice spacings $a_N\downarrow0$, one increasing sequence
of physical periods $P_p\uparrow\infty$, bare Wilson coefficients
$\beta_{N,p}$, and local observables \eqref{eq:headline-insertion} satisfying
\eqref{eq:headline-jet-normalization}, all on a single Wilson construction.
The coefficients satisfy \eqref{ex-running-input}. These choices precede
any choice of a subsequence of periods.

For every pair of strictly increasing index maps $(u,p)$ there is a
refinement $(u\circ v,p\circ q)$ and one common system of limiting fields,
correlations and phases with the following properties.
\begin{enumerate}
\item \emph{Wilson limits and field labels.}
The physical fields of degree at most $d$ are quotients
$\mathcal F_d=V_d/\mathcal N_d$, where $\mathcal N_d$ consists of labels
whose insertion vanishes in all separated mixed correlations with arbitrary
spectators. They have compatible injective grade inclusions, a unit,
conjugation, derivatives and spacetime symmetry actions. Write $[v]_d$ for
the field represented by $v\in V_d$. At grade $d$ the unit is represented
by $a_{r_d}^{d}\mathbf1$, not in general by the constant polynomial
$\mathbf1$. The grade inclusions are the compatible maps constructed in
Subsection~\ref{ex-phase-fields}; they need not send the class of a polynomial
to the class of the same polynomial at a larger grade. For every finite
list of graded labels $v_i\in V_{d_i}$, $1\le i\le n$, and
every smooth compactly supported test $f$ away from the coincidence diagonals,
\begin{equation}\label{eq:headline-wilson-limit}
\begin{split}
 \lim_{k\to\infty}\lim_{N\to\infty}
 a_N^{4n}\!\sum_{\boldsymbol x} f(\boldsymbol x)
 \mathbb E_{\mu_{N,p(q(k))}}
 \!\left[\prod_{i=1}^n I^{(d_i)}_{N,p(q(k))}(v_i;x_i;U)\right]\\
 =\int f(\boldsymbol x)S^{\rm mix}_{[v_1]_{d_1},\ldots,[v_n]_{d_n}}
                     (\boldsymbol x)\,d\boldsymbol x .
\end{split}
\end{equation}
Here $\mu_{N,p}$ is the Wilson probability measure of
Subsection~\ref{ex-lattice}, and the sum uses the physical lattice points in
its fundamental box, with the convention of \eqref{ex-smearing}.
The inner limit exists along the \emph{full} mesh sequence at every
suitable fixed period. One refinement works for all labels, arities and
tests. Auxiliary field comparisons may use the refined mesh $u\circ v$;
the Wilson limit displayed here does not.

\item \emph{Phases of that same limit.}
There is a compact metrizable space $X$ with a full-support probability
measure $\nu$ and smooth separated correlations $S^\chi$, multilinear in
the fields, such that
\[
 S^{\rm mix}_{\boldsymbol A}(\boldsymbol x)
   =\int_X S^\chi_{\boldsymbol A}(\boldsymbol x)\,d\nu(\chi).
\]
These are the characters of the commutative algebra obtained by compressing
the mixed theory's slab observables to its zero-energy space; $\nu$
represents its vacuum state. Thus the phase space and component correlations
come from the mixed Wilson limit, as in Sections~\ref{ex-phases}
and~\ref{ex-phase-selection}. Every component has reflection positivity
and the same field null relations as the mixed hierarchy; see
Sections~\ref{ex-phase-regularity} and~\ref{ex-phase-fields}.

\item \emph{Relativistic reconstruction.}
For every $\chi\in X$, these correlations reconstruct a nontrivial local
relativistic field theory on a Hilbert space with a common dense invariant
domain, cyclic vacuum, adjoints, tempered fields, strongly continuous
Poincar\'e symmetry and translation spectrum in the closed forward cone.
The translation-invariant vectors form precisely the vacuum line, and
nonzero field labels act nontrivially. Its vacuum distributions are the
analytic boundary values of $S^\chi$, in the sense of
Definition~\ref{ex-reconstructed} with $\chi$ in place of $\chi_*$.
The boundary distributions also obey the phase-mixture identity on joint
Schwartz tests. The analytic and field constructions are described in
Theorem~\ref{thm:joint-boundary} and Subsection~\ref{sec:word-reconstruction}.
There is a character $\chi_*\in X$, with one-dimensional ground space
as in Theorem~\ref{ex-pure-ground}, for which all the conclusions below
hold; they use these same data.

\item \emph{Local products and the physical stress tensor.}
There is one operator-product expansion for the mixed state and every
component, with the coefficient and remainder bounds
\eqref{ex-ope-coefficient}--\eqref{ex-ope-remainder}. The physical field
system contains a symmetric conserved stress tensor $T_{\mu\nu}$ whose
translation and rotation charges implement the field derivatives and
rotation action. Its normalization is that of
Subsection~\ref{sec:actual-stress-normalization}; its separated Ward identities
are those of Subsection~\ref{sec:physical-ward-propagation}.
In the selected theory set
\[
 \Theta=\sum_{\mu=1}^4T_{\mu\mu}
       -\left\langle\sum_{\mu=1}^4T_{\mu\mu}\right\rangle_{\chi_*}\mathbf1 .
\]
This very field, rather than a separately chosen grade-four field, satisfies
\[
 \lim_{r\downarrow0}r^8\log^2(r^{-1})
 S^{\chi_*}_{\Theta,\Theta}(-re_t/2,re_t/2)
       =\frac{3\dim_{\mathbb R}\mathfrak g}{\pi^4},
 \qquad e_t=(0,0,0,1).
\]
Its vacuum one-point functional vanishes and its two-point function has
the finite mass-shell-profile exclusion of
Subsection~\ref{sec:trace-profiles}. The comparison and limit use this same
normalized stress tensor; see Lemma~\ref{lem:trace-normalized-stress}
and Theorem~\ref{thm:trace-asymptotic}.
\end{enumerate}
The quotients, phases and reconstructed theories may depend on the
requested order and refinement. The lattice family does not.
\end{theorem}

\paragraph{Scope.}
This formulation selects conclusions of the construction; it does not
restate every source result. Contact extensions, nested operator products,
effective couplings and Maxwell scaling comparisons require separate precise
statements and are not asserted here. The running slope in
\eqref{ex-running-input} is $11C_{\mathfrak g}/(24\pi^2)$, with the
Casimir for the normalized trace inner product. Proposition~\ref{prop:shoot-public-running}
and \eqref{eq:public-family-running} identify this value with the bare
coefficients of this very Wilson family.

No uniqueness of the infinite-volume limit or independence of the spacing
sequence, local sections or blocking scheme is claimed. The result constructs
correlations and fields, not a probability measure on continuum connections,
Wilson-loop observables, particle states or a mass gap. The operator-product
bounds do not assert agreement with a perturbative series, and the trace test
does not exclude every generalized free-field description. These distinctions
matter in comparing the result with the Yang--Mills existence and mass-gap
problem; see \cite[Sections~1--2, 7--9]{YMG26}. In that comparison, the
axiomatic field theory and its curvature-polynomial representatives address
the first two requirements. The stress tensor, running law, operator-product
bounds and trace asymptotic give specific short-distance conclusions, but
not the full agreement with perturbative renormalization theory requested
in the third requirement. The fourth requirement, a positive spectral gap
above the vacuum, remains outside the result. The inverse reference maps
are constructed with quantitative bounds; explicit formulas for all their
matrix entries are not part of the assertion.

Forgetting the normalization, the character-algebra construction, the
reconstructions away from $\chi_*$ and the stress-tensor identification
gives, together with the unit, covering and token compatibility specified
below and proved in the field construction, the weaker public target.
Its insufficiency is discussed in
Subsection~\ref{ex-target-free-witness}. Theorem~\ref{thm:ex-construction-assembly} proves the stronger
same-data assertion: normalized insertions, field quotients, phase algebra,
every-character reconstruction and physical stress all use the same
extraction. Proposition~\ref{ex-conditional-existence} records the logical
projection to the weaker target. In particular the every-character
assertion in item~3 is supplied by its own vacuum argument, whereas the
trace asymptotic in item~4 concerns the distinguished selected character.
